\subsection{Dual of a Semisimple Algebra}

Let \(\Theta^{\mathrm{ss}}(\mathrm{A})\) be the socle of the left A-module \(\Theta (\mathrm{A})\) , that is, (VIII, p.~65) the largest semisimple submodule of \(\Theta (\mathrm{A})\) . We denote by \(\mathcal{S}_{\mathrm{K}}\) the set of classes of simple (left) A-modules that are finite-dimensional over K. When A is a semisimple algebra of finite degree over K, we have \(\mathrm{A}^{*} = \Theta (\mathrm{A}) = \Theta^{\mathrm{ss}}(\mathrm{A})\) because every left A-module is semisimple (VIII, p.~138, Proposition 4).

THEOREM 1. --- a) The set \(\Theta^{\mathrm{ss}}(\mathrm{A})\) consists of the coefficients of the finite-dimensional semisimple representations of A. It is an (A,A)-sub-bimodule of \(\mathrm{A}^*\) .

\begin{enumerate}
\def\labelenumi{\alph{enumi})}
\setcounter{enumi}{1}
\item
  For every \(S \in \mathcal{S}_{\mathrm{K}}\) , the isotypical component of \(\Theta(A)\) of type S is equal to \(\Theta_{\mathrm{S}}(\mathrm{A})\) . The left A-module \(\Theta^{\mathrm{ss}}(\mathrm{A})\) is the direct sum of the submodules \(\Theta_{\mathrm{S}}(\mathrm{A})\) , where S runs through \(\mathcal{S}_{\mathrm{K}}\) .
\item
  For every S in \(\mathcal{S}_{\mathrm{K}}\) , the right A-module \(S^{*}\) is simple, and \(\Theta_{\mathrm{S}}(\mathrm{A})\) is the isotypical component of type \(S^{*}\) of the right A-module \(\Theta(\mathrm{A})\) . The mapping that sends S to \(\operatorname{cl}(\mathrm{S}^{*})\) is a bijection from \(\mathcal{S}_{\mathrm{K}}\) to the set of classes of simple \(A^{o}\) -modules of finite dimension over K.
\item
  Viewed as a right A-module, \(\Theta (\mathrm{A})\) has socle \(\Theta^{\mathrm{ss}}(\mathrm{A})\) .
\end{enumerate}

Let E be a semisimple A-module of finite dimension over K. Every coefficient of E belongs to the image of an A-linear mapping from E to \(A^{*}\) (VIII, p.~378, Proposition 2) and therefore belongs to \(\Theta^{\mathrm{ss}}(\mathrm{A})\) . Conversely, let \(f \in \Theta^{\mathrm{ss}}(\mathrm{A})\) . Then the A-module Af is finite-dimensional over K and is semisimple. By the remark in VIII, p.~367, f is a coefficient of Af.

For every S in \(\mathcal{S}_{\mathrm{K}}\) , the isotypical component of \(\Theta(\mathrm{A})\) of type S is generated by the images of the A-linear mappings from S to \(A^{*}\) ; it is therefore equal to \(\Theta_{\mathrm{S}}(\mathrm{A})\) by Proposition 2, a) of VIII, p.~378. On the other hand, if S is a simple A-module that is not finite-dimensional over K, then the isotypical component of \(\Theta(\mathrm{A})\) of type S is zero: indeed, every simple submodule of \(\Theta(\mathrm{A})\) is monogenous and therefore finite-dimensional over K. Since the socle of \(\Theta(\mathrm{A})\) is the direct sum of its isotypical components (VIII, p.~65, Proposition 4), this proves b).

Let S be a simple A-module of finite dimension over K. Since the K-vector space S is not reduced to 0, the same holds for \(S^{*}\) . Let E be a submodule of the right A-module \(S^{*}\) ; its orthogonal \(E'\) in S is an A-submodule of S. Since S is simple, we have either \(E' = 0\) , in which case \(E = S^{*}\) , or \(E' = S\) , in which case E = 0. Hence \(S^{*}\) is a simple right A-module.

We have \(\Theta (\mathrm{A}^{\circ}) = \Theta (\mathrm{A})\) (VIII, p.~376); we identify right A-modules with left \(\mathrm{A}^{\circ}\) -modules. Since every vector space of finite dimension over K is isomorphic to its bidual, the above proves that the mapping \(\mathrm{S}\mapsto \mathrm{cl}(\mathrm{S}^{*})\) is a bijection from \(\mathcal{S}_{\mathrm{K}}\) to the set of classes of simple \(\mathrm{A}^{\circ}\) -modules of finite dimension over K. Now, for S in \(\mathcal{S}_{\mathrm{K}}\) , the isotypical component of \(\Theta (\mathrm{A}^{\circ})\) of type \(\mathrm{S}^*\) is equal to \(\Theta_{\mathrm{S}^*}(\mathrm{A}^{\circ})\) by assertion b) applied to the algebra \(\mathrm{A}^{\circ}\) , and we have \(\Theta_{\mathrm{S}^*}(\mathrm{A}^{\circ}) = \Theta_{\mathrm{S}}(\mathrm{A})\) . Assertions c) and d) follow immediately.

COROLLARY 1. --- Every simple A-module of finite dimension over K is isomorphic to an A-submodule of \(\Theta(A)\) .

This follows from Theorem 1, b) because we have \(\Theta_{\mathrm{E}}(\mathrm{A}) \neq 0\) for every nonzero A-module E.

COROLLARY 2. --- If two simple A-modules of finite dimension over K have a common nonzero coefficient, then they are isomorphic.

Let \(S_{1}\) and \(S_{2}\) be simple A-modules of finite dimension over K with a common nonzero coefficient. We then have \(\Theta_{S_{1}}(A) \cap \Theta_{S_{2}}(A) \neq 0\) . Now, \(\Theta_{S_{1}}(A)\) is isotypical of type \(S_{1}\) , and \(\Theta_{S_{2}}(A)\) is isotypical of type \(S_{2}\) . So \(S_{1}\) is isomorphic to \(S_{2}\) .

By Theorem 1, \(\Theta^{\mathrm{ss}}(\mathrm{A})\) is an (A,A)-sub-bimodule of \(\mathrm{A}^*\) , the direct sum of the (A,A)-bimodules \(\Theta_{\mathrm{S}}(\mathrm{A})\) for S running through \(\mathcal{S}_{\mathrm{K}}\) . These bimodules are pairwise nonisomorphic because they are already nonisomorphic as left A-modules.

Fix an S in \(\mathcal{S}_{\mathrm{K}}\) . Denote by D the opposite field of \(\operatorname{End}_{\mathrm{A}}(\mathrm{S})\) , and view S as a right D-module and \(S^{*}\) as a left D-module. With these conventions, S is an (A, D)-bimodule, \(S^{*}\) is a (D, A)-bimodule, and \(S \otimes_{D} S^{*}\) is an (A, A)-bimodule.

PROPOSITION 4. --- a) There exists a group homomorphism \(\lambda_{\mathrm{S}}\) from \(\mathrm{S} \otimes_{\mathrm{D}} \mathrm{S}^{*}\) to \(\Theta_{\mathrm{S}}(\mathrm{A})\) , characterized by

\[
\lambda_ {\mathrm{S}} (x \otimes x ^ {*}) = c _ {\mathrm{S}} (x, x ^ {*})\tag{6}
\]

for \(x \in S\) and \(x^{*} \in S^{*}\) . This mapping is an isomorphism of (A, A)-bimodules.

\begin{enumerate}
\def\labelenumi{\alph{enumi})}
\setcounter{enumi}{1}
\tightlist
\item
  The (A, A)-bimodule \(\Theta_{\mathrm{S}}(\mathrm{A})\) is simple.
\end{enumerate}

The mapping \(\gamma_{\mathrm{S}}: \mathrm{S} \otimes_{\mathrm{K}} \mathrm{S}^{*} \to \Theta_{\mathrm{S}}(\mathrm{A})\) characterized by the formula \(\gamma_{\mathrm{S}}(x \otimes x^{*}) = c_{\mathrm{S}}(x, x^{*})\) is (A, A)-linear and surjective and satisfies \(\gamma_{\mathrm{S}}(xd \otimes x^{*}) = \gamma_{\mathrm{S}}(x \otimes dx^{*})\) for \(x \in \mathrm{S}\) , \(x^{*} \in \mathrm{S}^{*}\) , and \(d \in \mathrm{D}\) . It therefore defines a surjective (A, A)-linear mapping \(\lambda_{\mathrm{S}}: \mathrm{S} \otimes_{\mathrm{D}} \mathrm{S}^{*} \to \Theta_{\mathrm{S}}(\mathrm{A})\) characterized by formula (6).

Let us prove that \(S \otimes_{D} S^{*}\) is a simple (A, A)-bimodule. By Corollary 2 of VIII, p.~63, every (A, A)-sub-bimodule of \(S \otimes_{D} S^{*}\) is of the form \(S \otimes_{D} H\) , where H is a (D, A)-sub-bimodule of \(S^{*}\) . Since \(S^{*}\) is a simple right A-module (VIII, p.~380, Theorem 1, c)), we have H = 0 or \(H = S^{*}\) ; the assertion follows.

The homomorphism \(\lambda_{S}\) is (A, A)-linear and nonzero; since \(S \otimes_{D} S^{*}\) is a simple bimodule, \(\lambda_{S}\) is injective (VIII, p.~47, Proposition 2, a)).

Remark. --- When the field K is algebraically closed, we have D = K by Theorem 1 of VIII, p.~47, and \(\lambda_{S}\) is an isomorphism of (A, A)-bimodules from \(S \otimes_{K} S^{*}\) to \(\Theta_{\mathrm{S}}(\mathrm{A})\) .
% semantic-fragments: legacy-input-seam
\relax{}
